% !TEX program = xelatex
% AIPL Overview — bird's-eye index of the project
% Generated from AIPL_OVERVIEW.md, 2026-05-18.

\documentclass[11pt,a4paper]{article}

\usepackage[a4paper,margin=22mm]{geometry}
\usepackage{xeCJK}
\usepackage{fontspec}
\usepackage{amsmath,amssymb}
\usepackage{listings}
\usepackage{xcolor}
\usepackage{booktabs}
\usepackage{tabularx}
\usepackage{longtable}
\usepackage{array}
\usepackage{hyperref}
\usepackage{enumitem}

\setCJKmainfont{Hiragino Mincho ProN}
\setCJKsansfont{Hiragino Sans}
\setCJKmonofont{Hiragino Sans}
\setmonofont{Menlo}[Scale=0.85]

\hypersetup{
  colorlinks=true,
  linkcolor=blue!50!black,
  urlcolor=blue!60!black,
  citecolor=black,
  pdftitle={AIPL Overview},
  pdfauthor={Yasushi Kodama}
}

\definecolor{codebg}{RGB}{248,248,242}
\definecolor{codecmt}{RGB}{120,120,120}
\definecolor{codekw}{RGB}{30,80,160}
\definecolor{codestr}{RGB}{160,40,40}

\lstdefinelanguage{aipl}{
  keywords={class, method, function, return, var, float, int, new, send, call, become,
            if, else, while, do, self, sender, now, future, await, true, false, pub, linear,
            scope, where, and, or, not},
  morekeywords={[2]match, with, otherwise, where},
  sensitive=true,
  comment=[l]{//},
  morestring=[b]",
  showstringspaces=false,
}

\lstset{
  basicstyle=\ttfamily\footnotesize,
  backgroundcolor=\color{codebg},
  keywordstyle=\color{codekw}\bfseries,
  commentstyle=\color{codecmt}\itshape,
  stringstyle=\color{codestr},
  frame=single,
  framesep=4pt,
  rulecolor=\color{codebg!60!black},
  breaklines=true,
  columns=fullflexible,
  keepspaces=true,
}

\title{\textbf{AIPL Overview}\\
\large Actor-based Intelligent Parallel Language --- bird's-eye index}
\author{Yasushi Kodama}
\date{2026-05-18}

\begin{document}
\maketitle

\begin{abstract}
\noindent AIPL は以下の 3 要素を組合せた研究プラットフォーム:
(1) マルチランタイム actor 言語 (OCaml + Python + Browser-JS),
(2) MAP-Elites 進化計算による次世代言語予測器 (\texttt{aice-evolution-v2}),
(3) AIPL 自身による型検査・評価系の self-hosting (\texttt{aipl-self-host}).
本ドキュメントはプロジェクト全体の俯瞰索引であり,各アーティファクトと
それらの繋がりを名指しで示す.
\end{abstract}

\tableofcontents
\vspace{8pt}

% ============================================================
\section{プロジェクト構成}
% ============================================================

AIPL は 3 つの柱で構成される:

\begin{enumerate}\setlength{\itemsep}{2pt}
\item \textbf{マルチランタイム actor 言語}
   (OCaml ネイティブ / Python / Browser-JS).
\item \textbf{次世代言語予測器} (\texttt{aice-evolution-v2}): GA で
   言語設計の普遍軸を経験的に発見する.
\item \textbf{Self-host}: AIPL のパース・型検査・評価・スケジューラを
   AIPL 自身で実装したもの (\texttt{aipl-self-host/}).
\end{enumerate}

3 ランタイムは HTTP \texttt{/api/json/call} で相互運用可能で,
\texttt{docker/cross/} に Python $\leftrightarrow$ OCaml の cross-language
remote-actor デモが収録されている.

% ============================================================
\section{アーキテクチャ図}
% ============================================================

\begin{lstlisting}
[aice-evolution-v2 (next-gen prediction)]
   .aice spec  ---->  MAP-Elites GA  ---->  elite_map / lineage / report

[src/python-aipl (research runtime)]
   lexer (Lark grammar)  -->  parser -> AST  -->  aipl_typeck (Phase 11-16)
                                                        |
                                                        v
                                          aipl_interp (actor + scheduler + builtins)

[src/ (OCaml runtime)]
   lexer.mll  -->  parser.mly  -->  infer.ml (HM-style)  -->  eval_thread.ml

[aipl-self-host (AIPL in AIPL)]
   Level A: metacircular eval
     -->  B-1 typeck  -->  B-2 effects  -->  B-3 channels
     -->  B-4 linear  -->  B-5 owned
     -->  C lexer+parser  -->  C-2 actor scheduler  -->  C-3 now/Reply

[docker/cross (cross-language remote actors)]
   aipl-python:cross  <-->  aipl-ocaml:cross  via /api/json/call

[aipl_dist (AIPL v2 Distributed runtime, Phase E後)]
   AIPL_DIST_ENABLE=1 で 8 opt-in 機能を重ね合わせ
   既存 .aipl / .aipl は無改変
\end{lstlisting}

% ============================================================
\section{Phase 11--17: ランタイム機能}
% ============================================================

Python ランタイムの主要 Phase と対応サンプル:

\begin{center}
\begin{tabular}{lll}
\toprule
phase & feature & source / sample \\
\midrule
11a--e & typed annotations / unions / generics / length-arrays & \texttt{Typecheck*.aipl} \\
12 & capability effects \texttt{!\{fs,ai,net,mut\}} & \texttt{Effects.aipl} \\
13 & CSP channels & \texttt{Channels.aipl} \\
14 & linear / use-after-move & \texttt{Linear.aipl} \\
15 & owned (\texttt{pub} field visibility) & \texttt{Owned.aipl} \\
16 & transient cast at any-boundary & \texttt{Transient.aipl} \\
17 & structured concurrency (\texttt{scope \{ future ... \}}) & \texttt{Phase17\_StructuredConc.aipl} \\
\bottomrule
\end{tabular}
\end{center}

各 Phase に \texttt{\_test\_*.py} 単体テスト (\texttt{src/python-aipl/}) が付属.

% ============================================================
\section{Phase C--E2: HM 型推論 + 分散ランタイム (2026-05-17 以降)}
% ============================================================

Phase 11 系の gradual 静的型検査に加え,Python 版 AIPL は Phase C 以降で
\textbf{constraint-based Hindley--Milner 型推論 + Z3 refinement} を備えるよう
進化した.実装ファイル \texttt{src/python-aipl/aipl\_inference.py} (約 1255 LOC).

\begin{center}
\begin{tabular}{lll}
\toprule
phase & feature & 関連レポート \\
\midrule
C & 基本 HM + Z3 refinement (Int) & \texttt{PHASE\_C\_REPORT.md} \\
D-1 & Cross-class 推論 & \texttt{PHASE\_D\_REPORT.md} \\
E-$\alpha$ & \texttt{where} 句を AIPL 表層に追加 & \texttt{PHASE\_E\_REPORT.md} \\
E-$\beta$ & Actor field 共有 (method 横断) & \texttt{PHASE\_E\_BETA\_REPORT.md} \\
E-$\gamma$ & Record structural typing & \texttt{PHASE\_E\_GAMMA\_REPORT.md} \\
E-$\gamma$-R & Real/Rat refinement (Z3 Real 理論) & \texttt{PHASE\_E\_GAMMA\_R\_REPORT.md} \\
E-2 & typeck $\times$ inference 統合 CLI \texttt{--check} & \texttt{PHASE\_E\_2\_REPORT.md} \\
\bottomrule
\end{tabular}
\end{center}

これと並行して,進化計算 (Run 1 + Run 2) で発見された分散ランタイム設計を
\texttt{aipl\_dist.py} (591 LOC) として実装:

\begin{center}
\begin{tabular}{p{2.6cm}p{5.5cm}p{6cm}}
\toprule
機能 & env var & 説明 \\
\midrule
I-1 routing & \texttt{AIPL\_ROUTE} & actor 名 $\to$ tag テーブル \\
I-2 structured\_log & \texttt{AIPL\_DIST\_LOG\_FILE} & ND-JSON 1 行 / event \\
I-3 token\_budget & \texttt{AIPL\_DIST\_RPM} / \texttt{TPM} & 60 秒スライディング窓 \\
I-4 checkpoint & \texttt{AIPL\_DIST\_CHECKPOINT\_DIR} & atomic save/restore \\
IQ quarantine & \texttt{AIPL\_DIST\_QUARANTINE\_TTL} & 自動隔離 \\
IM-1 quorum & \texttt{AIPL\_DIST\_QUORUM\_PROVIDERS} & 並列 multi-provider \\
IM-2 subtree & \texttt{AIPL\_DIST\_SUBTREE\_QUARANTINE} & Erlang OTP blast-radius \\
\bottomrule
\end{tabular}
\end{center}

詳細は \texttt{USER\_MANUAL.md} §6 / §7 (PDF: \texttt{docs/AIPL\_User\_Manual.pdf}) を参照.

% ============================================================
\section{Self-host (\texttt{aipl-self-host/})}
% ============================================================

\begin{center}
\begin{tabular}{lll}
\toprule
Level & AIPL 自身が何をするか & smoke \\
\midrule
A & \texttt{eval(ast, env)} metacircular & 4/4 \\
B-1 & Phase 11 型検査器 & 6/6 \\
B-2 & Phase 12 capability-effects 検査器 & 4/4 \\
B-3 & Phase 13 channel element-type 検査器 & 5/5 \\
B-4 & Phase 14 linear / use-after-move & 4/4 \\
B-5 & Phase 15 owned (pub 可視性) & 4/4 \\
C & lexer + parser + eval pipeline & 3/3 \\
C-2 & actor scheduler + mailbox & 6/6 \\
C-3 & + \texttt{now} / Reply 同期返信 & 1/1 \\
\midrule
\multicolumn{2}{l}{\textbf{合計}} & \textbf{37/37 smoke} \\
\bottomrule
\end{tabular}
\end{center}

各 Level B-N の検査器は \texttt{SampleSelfConsistency.aipl} で
自身にも適用済 — phase monotonicity を実証.

% ============================================================
\section{aice-evolution-v2 --- 次世代言語予測}
% ============================================================

\texttt{aice-evolution-v2/} の MAP-Elites GA は \texttt{.aice} 仕様を入力に取り,
cross-task 評価 + LLM-directed mutation で言語設計の普遍軸を経験的に導出する.

\begin{center}
\begin{tabular}{ll}
\toprule
.aice & 予測内容 \\
\midrule
\texttt{NextLanguagePrediction.aice} & Phase 9 baseline (5 普遍軸) \\
\texttt{AIPLSelfHost\_A\_Metacircular.aice} 他 & self-host 設計空間 (Level 別) \\
\texttt{AIPLPostC2\_NextGen.aice} & Phase 17 候補 (post-Level-C-2 seed) \\
\texttt{AIPLPhase17\_Structured.aice} & Phase 17 structured-concurrency 精緻化 \\
\texttt{AIPL\_v2\_TypeInference.aice} & Phase C-E2 (型推論強化) \\
\texttt{AIPL\_v2\_Distributed.aice} & I0003/I0023/I0036 (ランタイム層分散化) \\
\bottomrule
\end{tabular}
\end{center}

\paragraph{Phase 9 予測 (実装済)}
5 軸すべて Phase 11--15 として実装:

\begin{lstlisting}
type_safety          = high / dependent     -> Phase 11
effect_handling      = capability           -> Phase 12
concurrency_model    = csp_channels         -> Phase 13
ownership_model      = linear               -> Phase 14
state_representation = symbol_owned         -> Phase 15
\end{lstlisting}

\paragraph{Phase 17 予測 (実装済)}
post-Level-C-2 seed で再走した時の trend top 3:

\begin{lstlisting}
concurrency_model       +0.50    <- 最強 (= structured)
type_safety             +0.25
ownership_model         +0.25
\end{lstlisting}

elite は \texttt{concurrency\_model = structured + ownership\_model = linear +
type\_safety = dependent} で,これがそのまま Phase 17 (\texttt{scope \{ future ... \}}
linear future) として実装された.

\paragraph{Phase E-2 後 (Distributed 設計探索)}
\texttt{AIPL\_v2\_Distributed.aice} を Run 1 + Run 2 (各 24 min,1400+ LLM calls)
で探索した結果:

\begin{itemize}\setlength{\itemsep}{0pt}
\item Run 2 throughput\_first 勝者: I0036 (Erlang OTP コア,
       \texttt{restart\_subtree + quorum\_replicate + multi\_process\_local}).
\item Run 2 hang\_resilience 勝者: I0023 (\texttt{quarantine\_and\_skip +
       cluster\_with\_scheduler}).
\item Run 2 balanced 勝者: I0003 (低コスト変種,
       \texttt{token\_budget\_aware + same\_provider\_retry}).
\end{itemize}

3 候補のうち I0003 + IQ + IM-1 + IM-2 が \texttt{aipl\_dist.py} として
実装済 (上記 Phase C-E2 セクション参照).

\paragraph{可視化}
\texttt{aice-evolution-v2/viz/} 下:
\texttt{index.html} (サマリ),
\texttt{level\_\{A,B1,B2,C\}.html} (Chart.js ダッシュボード),
\texttt{level\_\{A,B1,B2,C\}\_tree.html} (SVG 系統樹),
\texttt{postC2.html} / \texttt{phase17.html} (post-C2 + 精緻化 run).

% ============================================================
\section{型ソウンドネスの境界}
% ============================================================

\texttt{docs/AIPL\_Type\_Soundness\_Report.\{tex,pdf\}} は AIPL の gradual
型システムが何を健全に保証し,どの境界で実行時検査に gracefully
退化するかを論じる.Phase 16 で最優先 unsoundness gap (any $\to$ T 境界)
を transient cast で塞いだ.

% ============================================================
\section{クロス言語 remote actors}
% ============================================================

\texttt{docker/cross/} は 2 つの Docker image
(\texttt{aipl-python:cross} と \texttt{aipl-ocaml:cross}),
\texttt{docker-compose.yml},および fallback の \texttt{docker run} ドライバを
収録.両 image は同一の HTTP wire format (\texttt{POST /api/json/call}) を
話すので,Python actor から OCaml actor を呼べる (逆も同様).

% ============================================================
\section{推奨読書順}
% ============================================================

新規 reader に推奨する path:

\begin{enumerate}\setlength{\itemsep}{0pt}
\item \texttt{USER\_MANUAL.md} (PDF: \texttt{docs/AIPL\_User\_Manual.pdf}) ---
       言語リファレンス.§6 (Phase C--E2 型推論) と §7 (AIPL v2 Distributed
       Runtime / \texttt{aipl\_dist}) が post-Phase-11 の新規部分.
\item \texttt{docs/AIPL\_Type\_Soundness\_Report.pdf} --- 型システムの
       証明と未解決部.
\item \texttt{aipl-self-host/level-a/README.md} $\to$ \texttt{level-b\{1,2,3,4,5\}} ---
       AIPL 自身の型システム実装.
\item \texttt{aipl-self-host/level-c\{,2,3\}/README.md} --- パーサと
       スケジューラの self-host.
\item \texttt{aice-evolution-v2/README.md} $\to$ \texttt{docs/AIPL\_Phase17\_Prediction.md} ---
       次世代軸予測器の中身と現在の推奨.
\item \texttt{docker/cross/README.md} --- 言語間相互運用デモ.
\end{enumerate}

% ============================================================
\section{リポジトリ配置}
% ============================================================

\begin{lstlisting}
abclcp-project/
  AIPL_OVERVIEW.md            this overview (markdown)
  USER_MANUAL.md              language reference
  README.md                   build instructions
  Makefile                    OCaml + JS build targets
  docs/
    AIPL_Overview.{tex,pdf}                this file's PDF
    AIPL_User_Manual.{tex,pdf}             reference PDF
    AIPL_Type_Soundness_Report.{tex,pdf}
    AIPL_Phase17_Prediction.md
    ...
  src/
    *.ml / *.mll / *.mly      OCaml runtime
    python-aipl/              research runtime
      aipl_inference.py       Phase C-E2 HM + Z3
      aipl_dist.py            AIPL v2 Distributed (8 features)
      aipl_typeck.py          Phase 11-16 gradual checker
      aipl_interp.py          interp (+ I-5/I-6 hooks)
      aipl_runtime.py         actor runtime (+ IQ/IM hooks)
    browser-abcl/             browser-side runtime
  abclc/                      OCaml-side .aipl samples
  aice-evolution-v2/          next-gen-axis predictor
    examples/*.aice           GA specs
    src/                      GA implementation
    viz/                      HTML visualisations
  aice-pi-evolution/
    experiments/2026-05-17_aipl_v2_type_inference/
      PHASE_*_REPORT.md       all Phase C-E2 reports
      AIPL_v2_Distributed.{aice,ga.json}
      IMPL_I0003_MVP/         I0003 + IQ + IM impl, tests, samples
  aipl-self-host/             AIPL in AIPL (Levels A -> C-3)
  docker/cross/               cross-language Docker demo
\end{lstlisting}

% ============================================================
\section{現状サマリ}
% ============================================================

2026-05-18 時点で全体スモーク:

\begin{itemize}\setlength{\itemsep}{0pt}
\item OCaml runtime: 52/52 sample tests (REPL + aipl2c + cc/SDL2)
\item Browser JS: 7/7 syntax + 4/4 parse
\item Python runtime: 33/33 \texttt{.aipl} + 21/21 \texttt{.aipl} (Phase A-G)
\item Distributed cross-language: 8/8 (3-node mock)
\item Self-host: 37/37 (Level A $\to$ C-3)
\item \texttt{aipl\_dist} 単体テスト: 27/27
\item Distributed MVP デモ: 24/24 (8 feature $\times$ 3)
\end{itemize}

\textbf{合計 213+ テスト + デモがクリーン}.進化計算 (24 min $\times$ 2 run)
$\to$ I-MVP 実装 (710 LOC) $\to$ feature 別 24 sample で実演,という一連が
完走している.

\end{document}
